\documentclass[12pt]{article}

\usepackage{amsmath}
\usepackage{amssymb}
\usepackage{amsthm}
\usepackage{fullpage}

\newcommand{\R}{\mathbb{R}}
\newcommand{\Z}{\mathbb{Z}}
\newcommand{\Requests}{\mathcal{R}}
\newcommand{\T}{\mathcal{T}}
\newcommand{\I}{\mathcal{I}}
\newcommand{\E}{\mathcal{E}}
\renewcommand{\P}{\mathcal{P}}

\DeclareMathOperator*{\argmax}{argmax}
\DeclareMathOperator*{\argmin}{argmin}
\DeclareMathOperator*{\Exp}{\mathbf{E}}

\DeclareMathOperator*{\sign}{sign}
\DeclareMathOperator*{\maximize}{maximize}
\DeclareMathOperator*{\minimize}{minimize}
\DeclareMathOperator*{\subjectto}{subject\ to}

\newtheorem{theorem}{Theorem}
\newtheorem{lemma}[theorem]{Lemma}
\newtheorem{proposition}[theorem]{Proposition}
\newtheorem{conjecture}[theorem]{Conjecture}

\theoremstyle{definition}
\newtheorem{definition}[theorem]{Definition}

\begin{document}

\title{Pacing}
\author{D\'avid P\'al}
\date{July 29, 2026}

\maketitle

\begin{abstract}
We show convergence results for the pacing problem.
\end{abstract}

\section{Introduction}

We consider an online optimization problem where an algorithm processes a
sequence of \emph{requests} $R_1, R_2, \dots, R_T$ one by one. Each request is a
non-empty subset of $\R^{M+1}$, finite or infinite. In round $t$, the algorithm
needs to choose an \emph{action} $a_t \in R_t$. The goal of the algorithm is to
solve approximately the optimization problem
\begin{equation}
\label{equation:optimization-problem}
\begin{aligned}
\maximize_{a_1, a_2, \dots, a_T} \qquad & \frac{1}{T} \sum_{t=1}^T a_{t,M+1} \\
\subjectto \qquad & \frac{1}{T} \sum_{t=1}^T a_{t,i} \le B_i && \qquad \text{for all $i=1,2,\dots,M$} \\
& a_t \in R_t && \qquad \text{for all $t=1,2,\dots,T$} \\
\end{aligned}
\end{equation}
where $B_1, B_2, \dots, B_M$ are real numbers known ahead of time.

Problems of the form \eqref{equation:optimization-problem} arise in revenue
management. For example, in online marketing, the request $R_t$ corresponds to
an opportunity for an advertiser to show an advertisement to a user. The
elements of $R_t$ correspond to the bids the advertiser can submit in an
advertisment auction. The coordinates of the action $a_t$ are the business
metrics the advertiser wishes to optimize, such as the probability that the user
clicks on the advertisement, the expected amount the advertiser pays to the
publisher, and the revenue the advertiser collects from the user.

Another example of a problem of the form \eqref{equation:optimization-problem}
arises in ride hailing. The request $R_t$ corresponds to a user request for a
ride. The elements of $R_t$ correspond to pairs of prices the ride hailing
platform can offer to the rider and the driver for the ride. The coordinates of
the action $a_t$ are the business metrics the ride hailing platform wishes to
optimize, such as the probability trip happening (i.e. both rider and driver
accept the ride), the expected revenue collected from the rider, and expected
profit for the ride hailing platform.

In applications, the elements of $R_t$ are often predicted by machine learning
models just before the action $a_t$ needs to be chosen. These predictions are
only proxies for the business metrics. The actual business metrics are revealed
only after the action has been chosen, sometimes with a significant delay.
However, in order to keep our model simple, we ignore the difference between the
predicted and the actual business metrics.

Without any assumptions on the sequence of requests, solving the optimization
problem, even approximately, is information-theoretically impossible; see
Appendix~\ref{section:impossibility-result} for a proof. We therefore make
several assumptions on the sequence of requests. First, we assume the requests
are drawn i.i.d. from some distribution $D$. The distribution $D$ is unknown to
the algorithm and can be arbitrary. Second, we assume the requests are uniformly
bounded. For concreteness, we assume $R_1, R_2, \dots, R_T \subseteq [-1,
1]^{M+1}$. Third, we assume that there exists an oracle that given a request $R
\subseteq [-1,1]^{M+1}$ and any $\epsilon > 0$ finds an $\epsilon$-net of $R$ of
size at most $\left\lceil \frac{2}{\epsilon} \right\rceil^{M+1}$.

The uniform boundedness assumption ensures that an influence of any single
requests on the optimization objective and the left-hand sides of the inequality
constraints is not too large. The choice of interval $[-1,1]$ is arbitrary and
can be replaced by any other fixed interval. At a technical level, the uniform
boundedness assumption allows us to use standard probabilistic concentration
inequalities and uniform convergence results. In the online marketing and ride
hailing applications, this assumption is satisfied.

The i.i.d. assumption is not satisfied in the applications mentioned earlier.
Both in online marketing and ride hailing, the sequence of requests has seasonal
patterns at multiple time scales (daily, weekly, etc.). Nevertheless, the i.i.d.
assumption is a good approximation if the $T$ requests span long enough time
period.

The $\epsilon$-net oracle is necessary to apply uniform convergence results and
to design an algorithm with efficient running time.

Under these assumptions, we design an algorithm that solves the online
optimization problem approximately. The idea behind the algorithm is to find an
approximately optimal policy. Let $\Requests$ be the set of all non-empty
subsets of $[-1,1]^{M+1}$. A \emph{policy} $\pi$ is a function $\pi:\Requests
\to [-1,1]^{M+1}$ that maps a subset $R \in \Requests$ to an action $\pi(R)$
such that $\pi(R) \in R$. A \emph{randomized policy} is a distribution over policies.
We denote by $\P$ the set of all randomized policies. The algorithm's goal is to
find a randomized policy $P \in \P$ that is an approximate solution of the
optimization problem
\begin{equation}
\label{equation:distribution-policy-optimization}
\begin{aligned}
\maximize_{P \in \P} \qquad & \Exp_{\substack{R \sim D \\ \pi \sim P}} \left[ \pi(R)_{M+1} \right] \\
\subjectto \qquad & \Exp_{\substack{R \sim D \\ \pi \sim P}} \left[ \pi(R)_i \right] \le B_i && \qquad \text{for all $i=1,2,\dots,M$} \\
\end{aligned}
\end{equation}
where the request $R$ is drawn randomly from the distribution $D$ and,
independently, the policy $\pi$ is drawn randomly from the distribution $P$. The
expectations above are taken with respect to both of these random draws, which
we make explicit using the subscript notation $R \sim D$, $\pi \sim P$. Using
Hoeffding's inequality, it is easy to show that an optimal (or approximately
optimal) solution $\pi$ of the optimization
problem~\eqref{equation:distribution-policy-optimization} solves the original
online optimization problem approximately.

Since the distribution $D$ is unknown, we cannot solve the policy optimization
problem \eqref{equation:distribution-policy-optimization} directly. Instead, in
round $t$, we use the sample $R_1, R_2, \dots, R_t$ as a proxy for the
distribution $D$. We find a randomized policy $P_t$ that approximately solves
the policy optimization problem on the sample
\begin{equation}
\label{equation:sample-policy-optimization}
\begin{aligned}
\maximize_{P_t \in \P} \qquad & \frac{1}{t} \sum_{\tau=1}^t \Exp_{\pi \sim P_t} \left[ \pi(R_\tau)_{M+1} \right] \\
\subjectto \qquad & \frac{1}{t} \sum_{\tau=1}^t \Exp_{\pi \sim P_t} \left[ \pi(R_\tau)_i \right] \le B_i && \qquad \text{for all $i=1,2,\dots,M$} \\
\end{aligned}
\end{equation}
We then use the randomized policy $P_t$ to choose the action $a_t$.

We study the relationship between the optimization problems
\eqref{equation:distribution-policy-optimization} and
\eqref{equation:sample-policy-optimization}. We show that good solutions to both
problems lie in a much smaller class of policies, which we call Lagrangian
policies. We show that a Lagrangian policy that is a good solution for one of
the optimization problems is a good solution for the other. Lagrangian policies
naturally arise from the formulation of the optimization
problem~\eqref{equation:sample-policy-optimization} as a linear program. The
dual of the linear program plays a central role in our analysis.

\subsection{Approximate solutions}

The optimization problems \eqref{equation:optimization-problem},
\eqref{equation:distribution-policy-optimization}, and
\eqref{equation:sample-policy-optimization} cannot be solved exactly. We thus
need to define appropriate notions of approximate solutions for these problems.
For any $\epsilon \ge 0$, we define the notion of $\epsilon$-feasible and
$\epsilon$-optimal solutions for these problems.

We say that a sequence of actions $a_1, a_2, \dots, a_T$ is
\emph{$\epsilon$-feasible} for the optimization problem
\eqref{equation:optimization-problem} if
\begin{align*}
\frac{1}{T} \sum_{t=1}^T a_{t,i} & \le B_i + \epsilon & \qquad \text{for all $i=1,2,\dots,M$,} \\
a_t \in R_t && \qquad \text{for all $t=1,2,\dots,T$.}
\end{align*}
Note that while the inequality constraints are relaxed, the constraint $a_t \in
R_t$ remains unchanged. We say that a randomized policy $P$ is
\emph{$\epsilon$-feasible} for the optimization problem
\eqref{equation:distribution-policy-optimization} if
$$
\Exp_{\substack{R \sim D \\ \pi \sim P}} \left[ \pi(R)_i \right] \le B_i + \epsilon \qquad \text{for all $i=1,2,\dots,M$.}
$$
We say that a randomized policy $P_t$ is \emph{$\epsilon$-feasible}
for the optimization problem \eqref{equation:sample-policy-optimization} if
$$
\frac{1}{t} \sum_{\tau=1}^t \Exp_{\pi \sim P_t} \left[ \pi(R_\tau)_i \right] \le B_i + \epsilon \qquad \text{for all $i=1,2,\dots,M$.}
$$

We say that a sequence of actions $a_1, a_2, \dots, a_T$ is
\emph{$\epsilon$-optimal} for the optimization problem
\eqref{equation:optimization-problem} if
$$
\frac{1}{T} \sum_{t=1}^T a_{t,M+1} \ge V - \epsilon \: ,
$$
where $V$ is the optimal value of the optimization problem
\eqref{equation:optimization-problem}. We say that a randomized policy $P$ is
\emph{$\epsilon$-optimal} for the optimization problem
\eqref{equation:distribution-policy-optimization} if
$$
\Exp_{\substack{R \sim D \\ \pi \sim P}} \left[ \pi(R)_{M+1} \right] \ge V_D - \epsilon
$$
where $V_D$ is the optimal value of the optimization problem
\eqref{equation:distribution-policy-optimization}.
We say that a randomized
policy $P_t$ is \emph{$\epsilon$-optimal} for the optimization problem
\eqref{equation:sample-policy-optimization} if
$$
\frac{1}{t} \sum_{\tau=1}^t \Exp_{\pi \sim P_t} \left[ \pi(R_\tau)_{M+1} \right] \ge V_t - \epsilon
$$
where $V_t$ is the optimal value of the optimization problem
\eqref{equation:sample-policy-optimization}. We adopt the convention that $V$,
$V_D$, $V_t$ equals $-\infty$ if the optimization
problem~\eqref{equation:optimization-problem},
\eqref{equation:distribution-policy-optimization},
\eqref{equation:sample-policy-optimization} is infeasible, respectively.

We are interested in solutions that are both $\epsilon$-feasible and
$\epsilon$-optimal. We remark that such solutions might exist even if the
optimization problems are infeasible.

\subsection{Notation and definitions}

We use the following notation in the rest of the paper. For $i=1,2,\dots,N$, we
define $e_i$ to be the vector in $\R^N$ with $i$-th entry equal to $1$ and all
other entries equal to zero, that is,
$$
e_i = \underbrace{(0,\dots,0,1,0,\dots,0)}_{\text{$1$ is at position $i$}} \: .
$$
For any $p \in [1, \infty]$, we denote by $\|\cdot\|_p$ the $p$-norm on $\R^{M+1}$.
That is, if $x = (x_1, \dots, x_{M+1}) \in \R^{M+1}$, then
$$
\|x\|_p =
\begin{cases}
\left(|x_1|^p + \dots + |x_{M+1}|^p\right)^{1/p} & \text{if $p \in [1, \infty)$,}  \\
\max \{|x_1|, \dots, |x_{M+1}|\} & \text{if $p = \infty$.}
\end{cases}
$$
For any $c \in \R^{M+1}$, any $r \ge 0$ and any $p \in [1, \infty]$,
we define the ball of radius $r$ centered at $c$ in $p$-norm,
$$
B_p(c,r) = \{x \in \R^{M+1} ~:~ \|x - c\|_p \le r \}
$$

We abuse the notation and allow a request to be a tuple of actions instead of a
set of actions. We use this convention only if the set of actions in the request
is finite. We allow an action to be repeated multiple times in the tuple. This
way, requests can be tuples of the same length. If $R = (a_1, a_2, \dots, a_N)$
is an $N$-tuple, the set notation $a \in R$ means that $a = a_i$ for some $i =
1,2,\dots,N$. If all requests are tuples of the same length $N$, a policy is
function $\pi: [-1,1]^{(M+1) \times N} \to [-1,1]^{M+1}$ such that $\pi(R) \in
R$.

\section{Discretization}
\label{section:discretization}

Requests that are arbitrary subsets of $[-1,1]^{M+1}$ are difficult to handle.
We therefore discretize them using $\epsilon$-nets. Formally, an
\emph{$\epsilon$-net} of a set $R \subseteq \R^{M+1}$ with respect to a norm
$\|\cdot\|$ is a subset $R' \subseteq R$ such that for every $x \in R$, there
exists $x' \in R'$ such that $\|x - x'\| \le \epsilon$. While the notion of
$\epsilon$-net can be considered for any norm, we will exclusively consider
$\epsilon$-nets with respect to the $\|\cdot\|_\infty$-norm.

We show that the size of $\epsilon$-net of $R$ can be upper bounded by
$\left\lceil \frac{2}{\epsilon} \right\rceil^{M+1}$; see
Proposition~\ref{proposition:size-of-epsilon-net} in
Appendix~\ref{section:size-of-epsilon-net}. The size of the $\epsilon$-net will
become important later.

We assume that we have an oracle that given any subset $R \subseteq [-1,1]^{M+1}$
and $\epsilon > 0$ finds an $\epsilon$-net of $R$ of size at most $\left\lceil
\frac{2}{\epsilon} \right\rceil^{M+1}$. We denote by $\Psi(R,\epsilon)$ the
$\epsilon$-net that the oracle outputs for input $R$ and $\epsilon$. Using the oracle,
we reduce the optimization problem~\eqref{equation:optimization-problem}
to the optimization problem
\begin{equation}
\label{equation:optimization-problem-discretized}
\begin{aligned}
\maximize_{a_1, a_2, \dots, a_T} \qquad & \frac{1}{T} \sum_{t=1}^T a_{t,M+1} \\
\subjectto \qquad & \frac{1}{T} \sum_{t=1}^T a_{t,i} \le B_i + \epsilon && \qquad \text{for all $i=1,2,\dots,M$} \\
& a_t \in \Psi(R_t,\epsilon) && \qquad \text{for all $t=1,2,\dots,T$}
\end{aligned}
\end{equation}

The problem \eqref{equation:optimization-problem-discretized} is a special case
of the optimization problem \eqref{equation:optimization-problem} where $R_t$ is
replaced by $\Psi(R_t,\epsilon)$ and $B_i$ is replaced by $B_i + \epsilon$. The
problem is special in the sense that the requests $\Psi(R_t,\epsilon)$ have
finite cardinality, and, in particular, the cardinality of each request is
bounded by $\left\lceil \frac{2}{\epsilon} \right\rceil^{M+1}$.

The optimization problem \eqref{equation:optimization-problem-discretized} is
closely related to the optimization problem
\eqref{equation:optimization-problem}. Lemma~\ref{lemma:epsilon-net-reduction}
below shows that a sequence of actions that is approximately feasible and
approximately optimal for optimization
problem~\eqref{equation:optimization-problem-discretized} is also approximately
feasible and approximately optimal for the optimization
problem~\eqref{equation:optimization-problem}. We refer to the error introduced
by this reduction as the \emph{discretization error}.

We remark that one can construct dicsretized versions of the optimization
problems \eqref{equation:distribution-policy-optimization} and
\eqref{equation:sample-policy-optimization}. Similar reductions and
corresponding analogues of Lemma~\ref{lemma:epsilon-net-reduction} exist for
these problems as well.

\begin{lemma}[$\epsilon$-net reduction]
\label{lemma:epsilon-net-reduction}
Let $\epsilon, \epsilon', \epsilon''$ be positive real numbers. Let
$a_1, a_2, \dots, a_T$ be a sequence of actions that is $\epsilon'$-feasible and $\epsilon''$-optimal for
the optimization problem~\eqref{equation:optimization-problem-discretized}. Then,
$a_1, a_2, \dots, a_T$ is $(\epsilon + \epsilon')$-feasible and $(\epsilon + \epsilon'')$-optimal
for the optimization problem~\eqref{equation:optimization-problem}.
\end{lemma}

\begin{proof}
Let $V$ be the optimal value of the optimization problem
\eqref{equation:optimization-problem}. Let $\epsilon''' > 0$ be a positive real
number. Then, there exists a sequence of actions $x'_1, x'_2, \dots, x'_T$ such that
\begin{align*}
\frac{1}{T} \sum_{t=1}^T a'_{t,M+1} & \ge V - \epsilon''' \: , \\
\frac{1}{T} \sum_{t=1}^T a'_{t,i} & \le B_i & \qquad \text{for all $i=1,2,\dots,M$,} \\
a'_t & \in R_t & \qquad \text{for all $t=1,2,\dots,T$.}
\end{align*}
In other words, $a'_1, a'_2, \dots, a'_T$ is feasible and $\epsilon'''$-optimal
for the optimization problem \eqref{equation:optimization-problem}.

We construct a sequence of actions $a''_1, a''_2, \dots, a''_T$. We define
$x''_t$ be the element of $\Psi(R_t,\epsilon)$ that is closest to $x'_t$ in the
$\|\cdot\|_\infty$-norm. Since $\Psi(R_t,\epsilon)$ is an $\epsilon$-net of
$R_t$,
\begin{align*}
\|a''_t - a'_t\|_\infty \le \epsilon & \qquad \text{for all $t=1,2,\dots,T$.}
\end{align*}
Therefore,
\begin{align*}
\frac{1}{T} \sum_{t=1}^T a''_{t,M+1} & \ge V - \epsilon''' - \epsilon \: , \\
\frac{1}{T} \sum_{t=1}^T a''_{t,i} & \le B_i + \epsilon & \qquad \text{for all $i=1,2,\dots,M$,} \\
a_t'' & \in \Psi(R_t,\epsilon) & \qquad \text{for all $t=1,2,\dots,T$.}
\end{align*}
In other words, $a''_1, a''_2, \dots, a''_T$ is a feasible and $(\epsilon''' +
\epsilon)$-optimal for the optimization
problem~\eqref{equation:optimization-problem-discretized}. Therefore, the
optimal value of the optimization
problem~\eqref{equation:optimization-problem-discretized} is at least $V -
\epsilon''' - \epsilon$. Since $\epsilon'''$ is arbitrary, the optimal value of
the optimization problem \eqref{equation:optimization-problem-discretized} is at
least $V - \epsilon$.

By assumption $a_1, a_2, \dots, a_T$ is $\epsilon'$-feasible and $\epsilon''$-optimal for
optimization problem \eqref{equation:optimization-problem-discretized}. This implies that
\begin{align*}
\frac{1}{T} \sum_{t=1}^T a'_{t,M+1} & \ge V - \epsilon - \epsilon'' \: , \\
\frac{1}{T} \sum_{t=1}^T a'_{t,i} & \le B_i + \epsilon + \epsilon' & \qquad \text{for all $i=1,2,\dots,M$,} \\
a_t' & \in \Psi(R_t,\epsilon) & \qquad \text{for all $t=1,2,\dots,T$.}
\end{align*}
Since $\Psi(R_t,\epsilon) \subseteq R_t$, this means that $a_1, a_2, \dots, a_T$
is $(\epsilon + \epsilon')$-feasible and $(\epsilon + \epsilon'')$-optimal for
the optimization problem \eqref{equation:optimization-problem}.
\end{proof}

As it will turn out, a suitable value for $\epsilon$ is $\epsilon = \sqrt{M/T}$.
This makes the discretization error commensurate with other sources of error
that we encounter later. In the rest of the paper, we assume that the requests
are $N$-tuples where $N = \left\lceil \frac{2}{\epsilon}
\right\rceil^{M+1}$.

\section{Lagrangian policies}

We consider policies parameterized by $M+1$ real numbers $\lambda_1, \dots,
\lambda_M, \lambda_{M+1}$. As will become clear later, the numbers $\lambda_1,
\dots, \lambda_M$ can be interpreted as the dual variables associated with the
constraints of the optimization problems. However, for now, they are arbitrary
real numbers. For any $(M+1)$-dimensional vector $\lambda = (\lambda_1, \dots,
\lambda_M, \lambda_{M+1}) \in \R^{M+1}$ we define a policy $\pi:[-1,1]^{(M+1)
\times N} \to [-1,1]^{M+1}$ by
\begin{equation}
\label{equation:lagrangian-policy}
\pi_\lambda(R) = \argmax_{a \in R} \lambda_1 a_1 + \lambda_2 a_2 + \dots + \lambda_M a_M + \lambda_{M+1} a_{M+1} \: .
\end{equation}
The ties in the $\argmax$ can be broken according to any consistent
deterministic rule. For concreteness, we break the ties using the lexicographic
ordering on $\R^{M+1}$.\footnote{The lexicographic ordering on $\R^{M+1}$ is
defined as follows. Let $x = (x_1, \dots, x_{M+1})$ and $y = (y_1, \dots,
y_{M+1})$ be two vectors in $\R^{M+1}$. Then $x$ is less than $y$
lexicographically if $x_1 < y_1$ or if $x_1 = y_1$ and $x_2 < y_2$ or if $x_1 =
y_1$ and $x_2 = y_2$ and $x_3 < y_3$ or if $x_1 = y_1$ and $x_2 = y_2$ and $x_3
= y_3$ and $x_4 < y_4$ and so on.} We call a policy of this form a
\emph{Lagrangian policy}. The reason for this name will become clear once we
formulate the optimization problem as a linear program and consider its dual.

We make three remarks about the Lagrangian policies. First, observe that
Lagrangian policies are deterministic.

Second, note that the policy corresponding to the all-zero vector $\lambda = (0,
0, \dots, 0)$ chooses the first action for any request. This policy is
uninteresting. For technical reasons and notational convenience, we include it
in certain theorems and exclude it in others. However, as it will turn out,
there exists a Lagrangian policy $\pi_\lambda$ with $\lambda \neq (0, 0, \dots,
0)$ that is close to an optimal policy of the optimization problems
\eqref{equation:sample-policy-optimization} and
\eqref{equation:distribution-policy-optimization}.

Third, note that for any positive real number $c$ and any vector $\lambda \in
\R^{M+1}$, the policies $\pi_{c \lambda}$ and $\pi_\lambda$ are identical. Thus,
without loss of generality, we can assume that $\|\lambda\|_2 \le 1$. If the
all-zero vector is excluded, we can assume that $\|\lambda\|_2 = 1$. Depending
on the context, the Lagrangian policies are parameterized by $\lambda$ that lies
in either $\R^{M+1}$, $\R^{M+1} \setminus \{0\}$, the unit ball, or the unit
sphere. The advantage of the unit ball and the unit sphere is that they are
compact sets.

\section{Failure of deterministic policies}

In general, deterministic policies are not good solutions of the optimization
problems~\eqref{equation:distribution-policy-optimization}~and~\eqref{equation:sample-policy-optimization}.
In this section, we provide the simplest possible example that illustrates the
point.

Suppose $M=1$, $B_1 = \frac{1}{2}$, $T \ge 1$ is arbitrary, $D$ is such that with probability one,
$$
R = ((1, 1), (0, 0)) \qquad \text{for all $t=1,2,\dots,T$.}
$$
The optimal randomized policy chooses actions $(1, 1)$ and $(0, 0)$ with
probability $1/2$ each. Any deterministic policy is either not a feasible
solution or not an optimal solution. The same problem arises for the
optimization problem~\eqref{equation:sample-policy-optimization}.

On the other hand, as we will see later, for distributions with densities these
problems vanish. Namely, deterministic Lagrangian policies are optimal solutions
for the optimization problem~\eqref{equation:distribution-policy-optimization}
when the distribution $D$ has a probability density. Furthermore, if the
requests are sampled i.i.d. from a probability distribution with a density, a
deterministic Lagrangian policy is close to an optimal solution for the
optimization problem~\eqref{equation:sample-policy-optimization}.


\section{Distributions with a density}

As explained in Section~\ref{section:discretization}, we can assume that the
requests are $N$-tuples of actions. Thus a distribution $D$ is a probability
distribution over $[-1,1]^{(M+1) \times N}$. So far, we have avoided making any
assumption on the distribution $D$. It is more convenient to assume that $D$ has
a probability density. Formally, that means that $D$ has a probability density
with respect to the Lebesgue measure on $\R^{(M+1) \times N}$. The fundamental
reason for assuming the existence of a probability density is to avoid having
ties in the argmax in equation~\eqref{equation:lagrangian-policy} defining
Lagrangian policies. The ties cause problems, as can be seen from
Proposition~\ref{proposition:computing-primal-from-dual}.

If the distribution $D$ does not have a probability density, we can always
create a ``nearby'' distribution $D'$ that does. One way to do so is to add
small independent noise to each entry of each action in the request. Namely,
suppose $R \sim D$. Let $Z$ be an $N$-tuple of vectors in $\R^{M+1}$ where each
entry is an i.i.d. random variable drawn from the uniform distribution over
$[-1,1]$. We construct a random sample $R'$ from $D'$ as
\begin{equation}
\label{equation:modified-sample}
R' = (1-\alpha)R + \alpha Z
\end{equation}
where $\alpha \in (0,1)$ is a ``small'' positive number. Since the sampling
error is of the order\footnote{The notation $\widetilde{O}(\cdot)$ hides
logarithmic factors in $N,M,T$.} $\widetilde{O}(\sqrt{M/T})$ and the rounding
error is of the order $O(M/T)$, a suitable value for $\alpha$ should be of
similar or smaller order.

Equation~\eqref{equation:modified-sample} defines the joint probability
distribution for $(R,R')$. The marginal distribution of $R$ is $D$, and the
marginal distribution of $R'$ is $D'$. Clearly, $R'$ and $R$ are not
independent. In other words, equation~\eqref{equation:modified-sample}
provides a non-trivial coupling between $D$ and $D'$. This coupling will play a
central role in this section.

We reduce the problem of finding a near-optimal solution for the
optimization problem~\eqref{equation:distribution-policy-optimization} for distribution
$D$ to a modified optimization problem for distribution $D'$. The modified
optimization problem is
\begin{equation}
\label{equation:distribution-policy-optimization-modified}
\begin{aligned}
\maximize_{P' \in \P} \qquad & \Exp_{\substack{R' \sim D' \\ \pi' \sim P'}} \left[ \pi'(R')_{M+1} \right] \\
\subjectto \qquad & \Exp_{\substack{R' \sim D' \\ \pi' \sim P'}} \left[ \pi'(R')_i \right] \le B_i' && \qquad \text{for all $i=1,2,\dots,M$} \\
\end{aligned}
\end{equation}
where
$$
B_i' = B_i(1 - \alpha) \: .
$$

First, we show that if the optimization
problem~\eqref{equation:distribution-policy-optimization} is feasible, then the
optimization problem~\eqref{equation:distribution-policy-optimization-modified}
is feasible as well. Second, we show that the optimal values of the two
optimization problems are related. These two claims are proved in
Lemma~\ref{lemma:modified-problem-is-feasible}.

Third, we show that any $\epsilon$-optimal policy $\pi':\Requests \to \Delta_N$
for the optimization
problem~\eqref{equation:distribution-policy-optimization-modified} can be
converted to an $\frac{\epsilon + \alpha}{1-\alpha}$-optimal policy
$\widetilde{\pi}$ for the optimization
problem~\eqref{equation:distribution-policy-optimization}; see
Theorem~\ref{theorem:policy-reduction}. The policy $\widetilde{\pi}$ is defined
as
$$
\widetilde{\pi}(R) = \pi'((1-\alpha)R + \alpha Z) \: ,
$$
where $Z$ is a random $N \times (M+1)$ matrix with entries drawn i.i.d. from the
uniform distribution over $[-1,1]$, independent of $R$. The output
$\widetilde{\pi}(R)$ is a random vector in $\Delta_N$. So, $\widetilde{\pi}$ is
not a function from $\Requests$ to $\Delta_N$. Nevertheless,
$\widetilde{\pi}(R)$ can be computed efficiently, provided that $\pi'(R')$ can
be computed efficiently.

\begin{lemma}[Modified problem is feasible]
\label{lemma:modified-problem-is-feasible}
Let $D$ be a probability distribution over $\Requests$. Let $\alpha \in (0,1)$.
Let $R$ be a random sample from $D$. Let $Z$ be an independent random $N \times
(M+1)$ matrix with entries drawn i.i.d. from the uniform distribution over
$[-1,1]$. Let $R' = (1-\alpha)R + \alpha Z$. Let $D'$ be the probability
distribution of $R'$. Let $V_D$ and $V_{D'}$ be the optimal values of the optimization
problems~\eqref{equation:distribution-policy-optimization} and
\eqref{equation:distribution-policy-optimization-modified}, respectively. If the
optimization problem~\eqref{equation:distribution-policy-optimization} is
feasible, then the optimization
problem~\eqref{equation:distribution-policy-optimization-modified} is feasible
as well. Furthermore,
$$
V_{D'} \ge (1 - \alpha) V_D \: .
$$
\end{lemma}

\begin{proof}
Let $\pi:\Requests \to \Delta_N$ be a feasible policy for the optimization
problem~\eqref{equation:distribution-policy-optimization}. We define the random
variable
$$
\widehat{\pi}(R') = \pi(R) = \pi \left( \frac{R' - \alpha Z}{1-\alpha} \right) \: .
$$
Using this random variable, we define the policy $\pi':\Requests \to \Delta_N$ as
$$
\pi'(R') = \Exp[ \widehat{\pi}(R') ~|~ R' ] \: .
$$
We claim that
\begin{equation}
\label{equation:reduction}
\Exp_{R' \sim D'} \left[ \sum_{i=1}^N \pi'(R')_i \cdot R'_{i,j} \right] =  (1 - \alpha) \Exp \left[ \sum_{i=1}^N \pi(R)_i \cdot R_{i,j} \right] \qquad \text{for any $j=0,1,2,\dots,M$} \: .
\end{equation}
Indeed,
\begin{align*}
\Exp_{R' \sim D'} \left[ \sum_{i=1}^N \pi'(R')_i \cdot R'_{i,j} \right]
& = \Exp \left[ \sum_{i=1}^N \Exp[ \widehat{\pi}(R') ~|~ R' ]_i \cdot R'_{i,j} \right] \\
& = \Exp \left[ \sum_{i=1}^N \Exp[ \widehat{\pi}(R')_i ~|~ R' ] \cdot R'_{i,j} \right] \\
& = \Exp \left[ \sum_{i=1}^N \Exp[ \widehat{\pi}(R')_i \cdot R'_{i,j} ~|~ R' ] \right] \\
& = \Exp \left[ \sum_{i=1}^N \widehat{\pi}(R')_i \cdot R'_{i,j} \right] \\
& = \Exp \left[ \sum_{i=1}^N \pi(R)_i \cdot R'_{i,j} \right] \\
& = \Exp \left[ \sum_{i=1}^N \pi(R)_i \cdot \left( (1 - \alpha) R_{i,j} + \alpha Z_{i,j} \right) \right] \\
& = (1 - \alpha) \Exp \left[ \sum_{i=1}^N \pi(R)_i \cdot R_{i,j} \right]  + \alpha \Exp \left[ \sum_{i=1}^N \pi(R)_i \cdot Z_{i,j} \right] \\
& = (1 - \alpha) \Exp \left[ \sum_{i=1}^N \pi(R)_i \cdot R_{i,j} \right] \: .
\end{align*}
In the calculation above, we used that $\Exp[ Z_{i,j} ] = 0$ and that $R$ and
$Z$ are independent.

Equation~\eqref{equation:reduction} implies that $\pi'$ is a feasible policy for
the optimization
problem~\eqref{equation:distribution-policy-optimization-modified}. Indeed, for
any $j=1,2,\dots,M$,
\begin{align*}
\Exp_{R' \sim D'} \left[ \sum_{i=1}^N \pi'(R')_i \cdot R'_{i,j} \right]
=  (1 - \alpha) \Exp \left[ \sum_{i=1}^N \pi(R)_i \cdot R_{i,j} \right]
\le (1 - \alpha) B_j
= B_j' \: .
\end{align*}

Equation~\eqref{equation:reduction} also implies the inequality $V_{D'} \ge (1 -
\alpha) V_D$. For any feasible solution $\pi$ of the optimization
problem~\eqref{equation:distribution-policy-optimization} there exists a
feasible solution $\pi'$ of the optimization
problem~\eqref{equation:distribution-policy-optimization-modified} such that
$$
\Exp_{R' \sim D'} \left[ \sum_{i=1}^N \pi'(R')_i \cdot R'_{i,0} \right] =  (1 - \alpha) \Exp \left[ \sum_{i=1}^N \pi(R)_i \cdot R_{i,0} \right] \: .
$$
Therefore, taking the supremum over all feasible solutions $\pi$, we get that
$$
V_{D'} \ge (1 - \alpha) V_D \: .
$$
\end{proof}

\begin{theorem}[Policy reduction]
\label{theorem:policy-reduction}
Let $D$ be a probability distribution over $\Requests$. Let $\alpha \in (0,1)$.
Let $R$ be a random sample from $D$. Let $Z$ be an independent random $N \times
(M+1)$ matrix with entries drawn i.i.d. from the uniform distribution over
$[-1,1]$. Let $R' = (1-\alpha)R + \alpha Z$. Let $D'$ be the probability
distribution of $R'$. Let $\pi':\Requests \to \Delta_N$ be an $\epsilon$-optimal
policy for the optimization
problem~\eqref{equation:distribution-policy-optimization-modified}. Let
$$
\widetilde{\pi}(R) = \pi'((1-\alpha)R + \alpha Z) \: .
$$
The policy $\widetilde{\pi}$ is an $\frac{\epsilon + \alpha}{1-\alpha}$-optimal
policy for the optimization
problem~\eqref{equation:distribution-policy-optimization}.
\end{theorem}

\begin{proof}
Straightforward calculation shows that for all $j=0,1,2,\dots,M$,
\begin{align*}
\Exp \left[ \sum_{i=1}^N \widetilde{\pi}(R)_i \cdot R_{i,j} \right]
& = \Exp \left[ \sum_{i=1}^N \pi'((1-\alpha)R + \alpha Z)_i \cdot R_{i,j} \right] \\
& = \Exp \left[ \sum_{i=1}^N \pi'(R')_i \cdot R_{i,j} \right] \\
& = \Exp \left[ \sum_{i=1}^N \pi'(R')_i \cdot \frac{R'_{i,j} - \alpha Z_{i,j}}{1-\alpha} \right] \\
& = \frac{1}{1-\alpha} \Exp \left[ \sum_{i=1}^N \pi'(R')_i \cdot R'_{i,j} \right] -  \frac{\alpha}{1-\alpha} \Exp \left[ \sum_{i=1}^N \pi'(R')_i \cdot Z_{i,j} \right] \: .
\end{align*}
We use that $\Exp \left[ \sum_{i=1}^N \pi'(R')_i \cdot Z_{i,j} \right]$ lies in
the interval $[-1,1]$ and that $\pi'$ is an $\epsilon$-optimal policy
for the optimization
problem~\eqref{equation:distribution-policy-optimization-modified}. Let $V_D$ and
$V_{D'}$ be the optimal values of the optimization
problems~\eqref{equation:distribution-policy-optimization}
and~\eqref{equation:distribution-policy-optimization-modified}.
Lemma~\ref{lemma:modified-problem-is-feasible} implies that $V_{D'} \ge (1 - \alpha)
V_D$. Combining all these facts we get that
\begin{align*}
\Exp \left[ \sum_{i=1}^N \widetilde{\pi}(R)_i \cdot R_{i,0} \right]
& = \frac{1}{1-\alpha} \Exp \left[ \sum_{i=1}^N \pi'(R')_i \cdot R'_{i,0} \right] -  \frac{\alpha}{1-\alpha} \Exp \left[ \sum_{i=1}^N \pi'(R')_i \cdot Z_{i,0} \right] \\
& \ge \frac{1}{1-\alpha} (V_{D'} - \epsilon) - \frac{\alpha}{1-\alpha}   \\
& \ge V_D - \frac{\epsilon + \alpha}{1-\alpha} \: ,
\end{align*}
and, for any $j=1,2,\dots,M$,
\begin{align*}
\Exp \left[ \sum_{i=1}^N \widetilde{\pi}(R)_i \cdot R_{i,j} \right]
& = \frac{1}{1-\alpha} \Exp \left[ \sum_{i=1}^N \pi'(R')_i \cdot R'_{i,j} \right] -  \frac{\alpha}{1-\alpha} \Exp \left[ \sum_{i=1}^N \pi'(R')_i \cdot Z_{i,j} \right] \\
& \le \frac{1}{1-\alpha} (B'_j + \epsilon) + \frac{\alpha}{1-\alpha} \\
& = B_j + \frac{\epsilon + \alpha}{1-\alpha} \: .
\end{align*}
In other words, $\widetilde{\pi}$ is an $\frac{\epsilon + \alpha}{1-\alpha}$-optimal policy for the optimization
problem~\eqref{equation:distribution-policy-optimization}.
\end{proof}

\appendix

\section{Impossibility result for adversarial requests}
\label{section:impossibility-result}

TODO

\section{Size of $\epsilon$-net}
\label{section:size-of-epsilon-net}

\begin{proposition}[Size of $\epsilon$-net]
\label{proposition:size-of-epsilon-net}
Let $R \subseteq [-1,1]^{M+1}$. Let $\epsilon > 0$.
Then $R$ has a finite $\epsilon$-net with respect to the $\|\cdot\|_\infty$-norm of size at most
$$
\left\lceil \frac{2}{\epsilon} \right\rceil^{M+1} \: .
$$
\end{proposition}

\begin{proof}
We can cover the interval $[-1,1]$ with $N = \left\lceil 2/\epsilon \right\rceil$ intervals
$$
[v_1 - \epsilon/2, v_1 + \epsilon/2], [v_2 - \epsilon/2, v_2 + \epsilon/2], \dots, [v_N - \epsilon/2, v_N + \epsilon/2]
$$
with centers at $v_k = \frac{(2k-1)\epsilon}{2}  - 1$ for $k = 1, 2, \dots, N$
and lengths $\epsilon$. Let $V = \{v_k ~:~ k = 1, 2, \dots, N\}$ be their
centers.

We can cover the set $[-1,1]^{M+1}$ with balls $B_\infty(c, \epsilon/2)$ with
centers ranging over $c \in V^{M+1}$. Formally,
$$
[-1,1]^{M+1} \subseteq \bigcup_{c \in V^{M+1}} B_\infty(c,\epsilon/2) \: .
$$
Let $C$ be set of the centers of those balls that intersect with $R$.
Formally,
$$
C = \{ c \in V^{M+1} ~:~ R \cap B_\infty(c,\epsilon/2) \neq \emptyset \}
$$
Clearly, the balls $B_\infty(c, \epsilon/2)$ with centers ranging over $c \in C$ cover the set $R$. Formally,
\begin{equation}
\label{equation:coverage-1}
R \subseteq \bigcup_{c \in C} B_\infty(c,\epsilon/2) \: .
\end{equation}

For every center $c \in C$, let $\phi(c)$ be a point in the intersection $R \cap
B_\infty(c,\epsilon/2)$. Let
$$
R' = \{\phi(c) ~:~ c \in C\} \: .
$$
We claim that $R'$ is an $\epsilon$-net of $R$. Since $\phi(c) \in R$, $R'
\subseteq R$. It remains to verify that for any $x \in R$ there exists a point a
point in $R'$ within distance $\epsilon$ of $x$. Let $x \in R$ be arbitrary.
The condition \eqref{equation:coverage-1} implies that there exists $c \in C$
such that $x \in B_\infty(c,\epsilon/2)$. Then, the point $\phi(c) \in R'$
satisfies
$$
\|\phi(c) - x\|_\infty \le \|\phi(c) - c\|_\infty + \|\phi(c) - c\|_\infty \le \epsilon/2 + \epsilon/2 = \epsilon \: .
$$

Finally, $R'$ is finite and has size at most
$$
|R'| \le |C| \le |V|^{M+1} = N^{M+1} = \left\lceil \frac{2}{\epsilon} \right\rceil^{M+1} \: .
$$
\end{proof}

\nocite{*}
\bibliographystyle{plain}
\bibliography{../bibliography}

\end{document}
